\documentclass[12pt]{article}

\usepackage[utf8]{inputenc}
\usepackage{latexsym,amsfonts,amssymb,amsthm,amsmath}
\usepackage{tikz}
\usepackage{stanli}
\usetikzlibrary{shapes,calc,through,3d,patterns}
\usetikzlibrary{arrows.meta,decorations.pathreplacing,angles,quotes}
\usepackage{multicol}
\usepackage{geometry}
\usepackage{mathtools}
\usepackage{siunitx}
\usepackage{booktabs}

\setlength{\parindent}{0in}
\setlength{\oddsidemargin}{-0.25in}
\setlength{\textwidth}{7in}
\setlength{\textheight}{9.3in}
\setlength{\topmargin}{-1in}
\setlength{\headheight}{18pt}

\newcommand\blfootnote[1]{%
    \begingroup
    \renewcommand\thefootnote{}\footnote{#1}%
    \addtocounter{footnote}{-1}%
    \endgroup
}

\definecolor{darkred}{rgb}{0.6,0,0}

\newcommand{\bolddelta}{\boldsymbol\delta}
\newcommand{\boldxi}{\boldsymbol\xi}



\begin{document}

\noindent ME473 - Dynamic finite element analysis of structures\hfill EPFL\\
Stefano Burzio \hfill 2025\\
\vspace{0.3cm}

\hspace*{\fill}\textbf{\Large Problem set 10 - solutions}\hspace*{\fill}

\hrulefill

\subsection*{Problem 1}
\textbf{1. Unconstrained and constrained degrees of freedom.} The bar is clamped at node 1, hence the corresponding displacement is constrained to zero: $d_1(t) = 0$. The unknown degrees of freedom are the displacements along the longitudinal axis at the unconstrained nodes 2 and 3, which are collected in the vector
\[
    \mathbf q_f(t) = \begin{bmatrix} d_2(t) \\ d_3(t) \end{bmatrix}
\]
To proceed, the global stiffness and mass matrices are partitioned accordingly, yielding:
\[
    \mathbf{K}_f = \frac{EA}{3\ell}
    \begin{bmatrix}
        16 & -8 \\
        -8 & 7
    \end{bmatrix}, \quad
    \mathbf{M}_f = \frac{\rho A \ell}{6}
    \begin{bmatrix}
        4 & 0 \\
        0 & 1
    \end{bmatrix}, \quad\text{and}\quad
    \mathbf r_f(t) = \begin{bmatrix} 0 \\ \delta(t) \end{bmatrix}.
\]
The differential equations governing the forced response of the bar is then written, taking into account the essential boundary condition at the clamped end, as follows:
\[
    \frac{\rho A \ell}{6}
    \begin{bmatrix}
        4 & 0 \\
        0 & 1
    \end{bmatrix} \begin{bmatrix} \ddot{d}_2(t) \\ \ddot{d}_3(t) \end{bmatrix}
    + \frac{EA}{3\ell}
    \begin{bmatrix}
        16 & -8 \\
        -8 & 7
    \end{bmatrix} \begin{bmatrix} d_2(t) \\ d_3(t) \end{bmatrix}
    = \begin{bmatrix} 0 \\ \delta(t) \end{bmatrix}
\]
The structure is initially at rest, thus the initial conditions are:
\begin{align*}
    {d}_2(0) =  {d}_3(0)         & = 0, \\
    \dot{d}_2(0) =  \dot{d}_3(0) & = 0.
\end{align*}
\textbf{2. Exact solution.} To determine the dynamic behavior of the bar under the action of an external load, we first evaluate the solution of equation  without the right-hand side (free response), which will then allow us to apply the modal superposition method using the modal parameters obtained in the first step.

We solve the generalized eigenvalue problem : $\mathbf{K}_f p_i = \lambda_i \mathbf{M}_f p_i $, for $i=1,2$. To simplify computations let:
\[
    \mathbf{K}_f = K_0
    \begin{bmatrix}
        16 & -8 \\
        -8 & 7
    \end{bmatrix}, \quad
    \mathbf{M}_f = M_0
    \begin{bmatrix}
        4 & 0 \\
        0 & 1
    \end{bmatrix}
\]
with $K_0 = \frac{EA}{3\ell}$ and  $M_0 = \frac{\rho A \ell}{6}$. We solve: $\left( \mathbf{K}_f - \lambda \mathbf{M}_f \right) \mathbf p =\mathbf 0$. The characteristic polynomial is:
\[
    \det\left(
    \begin{bmatrix}
            16K_0 - 4M_0 \lambda & -8K_0              \\
            -8K_0                & 7K_0 - M_0 \lambda
        \end{bmatrix}
    \right) = 0
\]
This simplifies to a quadratic in \( \lambda \):
$$4M_0^2 \lambda^2 - 44K_0 M_0 \lambda + 48K_0^2 = 0,$$
which reduces to:
$$\lambda^2 - 22 \frac{E}{\rho\ell} \lambda + 48\frac{E^2}{\rho^2\ell^2} = 0.$$
Solving this yields the eigenvalues:
$$ \lambda_1=(11-\sqrt{73})\frac{E}{\rho\ell^2}=2.4560\frac{E}{\rho\ell^2} \qquad\text{and}\qquad \lambda_2=(11+\sqrt{73})\frac{E}{\rho\ell^2}=19.5440\frac{E}{\rho\ell^2}.$$
and corresponding natural frequencies:
$$ \omega_1=1.5672\sqrt{\frac{E}{\rho\ell^2}} \qquad\text{and}\qquad \omega_2=4.4209 \sqrt{\frac{E}{\rho\ell^2}}.$$
Substituting the eigenvalues into the generalized eigenvalue problem allows us to extract the mode shapes $\mathbf{p}_{1}$ and $\mathbf{p}_{2}$, which we then normalize with respect to the diagonal mass matrix $\mathbf{M}_f$.
\[
    \mathbf{p}_{1} = \frac{1}{\sqrt{\rho A \ell}}
    \begin{bmatrix}
        1.0066 \\
        1.3952
    \end{bmatrix}, \quad
    \mathbf{p}_{2} =\frac{1}{\sqrt{\rho A \ell}}
    \begin{bmatrix}
        -0.6976 \\
        2.0133
    \end{bmatrix}
\]
Let \( \mathbf{P} = \begin{bmatrix} \mathbf p_1 & \mathbf p_2 \end{bmatrix} \) be the modal matrix of eigenvectors normalized with respect to \( \mathbf{M}_f \):
\[
    \mathbf{P}^T \mathbf{M}_f \mathbf{P} = \mathbf{I}.
\]
The modal parameters are now known. The modal superposition technique can be employed to determine the time response of the bar to a given excitation. We introduce the change of variables:
\[
    \mathbf q_f(t) = \mathbf{P} \mathbf z_f(t),
\]
which transforms the governing system into a set of uncoupled equations:
$$\ddot{\mathbf z}_f(t) + \boldsymbol{\Lambda} \mathbf z_f(t) = \mathbf{P}^T \mathbf r_f(t)$$
where \( \boldsymbol{\Lambda} = \text{diag}(\omega_1^2, \omega_2^2) \). Since \( \mathbf r_f(t) = \begin{bmatrix} 0 \\ \delta(t) \end{bmatrix} \), we have:
\[
    \mathbf s_f(t)=\mathbf{P}^T \mathbf r_f(t) = \begin{bmatrix} p_{12} \\ p_{22} \end{bmatrix} \delta(t).
\]
The system decouples into two second-order scalar differential equations:
($i=1,2$):
$$\ddot{z}_i(t)+\omega^2_iz_i(t)  =p_{i2}\delta(t).$$
The exact solution, using Duhamel's integral or Laplace transforms, is:
\begin{align*}
    z_i(t) & =\frac{1}{\omega_{i}} \int_0^t p_{i2}\delta(t - \tau) \sin(\omega_{i} \tau) \, \mathrm{d}\tau \\
           & = \frac{1}{\omega_{i}} p_{i2} \sin(\omega_{i}t).
\end{align*}
Returning to the physical coordinates via \( \mathbf{q}_f(t) = \mathbf{P} \mathbf{z}_f(t) \), we obtain the full dynamic response:
$$ \mathbf{q}_f(t) =\mathbf{p}_{1}z_1(t)+\mathbf{p}_{2}z_2(t).$$
Explicitly, the displacements at nodes 2 and 3 are:
\begin{align*}
    d_2(t) & = \frac{p_{11} \, p_{12}}{\omega_1} \sin(\omega_1 t)+ \frac{p_{21} \, p_{22}}{\omega_2} \sin(\omega_2 t) =\frac{1}{A\sqrt{\rho E}}\big( 0.8958 \sin(1.5672\, t)
    - 0.3178 \sin(4.4209\ t)\big)                                                                                                                                            \\
    d_3(t) & = \frac{(p_{12})^2}{\omega_1} \sin(\omega_1 t) + \frac{(p_{22})^2}{\omega_2} \sin(\omega_2 t)= \frac{1}{A\sqrt{\rho E}}\big( 1.2431 \sin(1.5672\, t)
    + 0.9168 \sin(4.4209\, t)\big)
\end{align*}


\textbf{3. Average acceleration Newmark method.} Recall from the previous section that the modal equations are:
\begin{align*}
    \ddot{z}_1(t) +  2.4560 \, \frac{E}{\rho A \ell^2} \, z_1(t) = 1.3952 \, \frac{1}{\sqrt{\rho A \ell}} \, \delta(t), \\
    \ddot{z}_2(t) + 19.544 \, \frac{E}{\rho A \ell^2} \, z_2(t) = 2.0133 \, \frac{1}{\sqrt{\rho A \ell}} \, \delta(t).
\end{align*}
At \( t = 0 \), the bar is at rest $z_1(0) = \dot{z}_1(0) = 0$ and $z_2(0) = \dot{z}_2(0) = 0$, thus $z_1^{(0)} =\dot{z}_1^{(0)}=0$ and $z_2^{(0)}  =\dot{z}_2^{(0)}=0.$
Based on the geometric and material properties of the bar and the definition of the time excitation function, the problem reduces to step-by-step iteration using the following dynamic relations:
\begin{equation}
    \begin{aligned}
        \ddot{z}_1^{(0)} + (8105.7)^2 z_1^{(0)}  & = 15747, \\
        \ddot{z}_2^{(0)} + (22866.2)^2 z_2^{(0)} & = 22723.
    \end{aligned}
    \label{eq:dyn}
\end{equation}
and for every time step index $k\geq 1$:
\begin{equation}
    \begin{aligned}
        \ddot{z}_1^{(k)} + (8105.7)^2 z_1^{(k)}  & = 0, \\
        \ddot{z}_2^{(k)} + (22866.2)^2 z_2^{(k)} & = 0.
    \end{aligned}
    \label{eq:dyn1}
\end{equation}
We begin by completing the initial phase of the Newmark's scheme: that is computing the accelerations at the initial time step using dynamic equations \eqref{eq:dyn}. Since the initial modal coordinates are zero, the initial modal accelerations become:
\begin{align*}
    \ddot{z}_1^{(0)} & = 15747 - (8105.7)^2 z_1^{(0)} = 15747  \\
    \ddot{z}_2^{(0)} & = 22723 - (22866.2)^2 z_2^{(0)} = 22723
\end{align*}

To compute the modal displacements, velocities, and accelerations at $t = \Delta t$, we employ the two kinematic relations associated with the Newmark average acceleration method, which for each $i = 1, 2$ take the form:
\begin{equation}
    \begin{aligned}
        z_i^{(k)}       & = \underbrace{z_i^{(k-1)} + \Delta t \, \dot{z}_i^{(k-1)} + 0.25 \Delta t^2 \ddot{z}_i^{(k-1)} }_\text{$\hat{z}_i^{(k-1)}$}+ 0.25 \Delta t^2\ddot{z}_i^{(k)} \\[10pt]
        \dot{z}_i^{(k)} & = \underbrace{\dot{z}_i^{(k-1)} + 0.5 \Delta t \ddot{z}_i^{(k-1)}}_\text{$\hat{\dot{z}}_i^{(k-1)}$} + 0.5 \Delta t \ddot{z}_i^{(k)}                          \\
    \end{aligned}
    \label{eq:kin}
\end{equation}
here $k = 1, 2, \dots$ denotes the time step index. It is clear that due to the implicit nature of the resolution method, the modal accelerations from the dynamic equations \eqref{eq:dyn1}, can only be extracted after being inserted into the kinematic relations \eqref{eq:kin}.

At the first iteration of the process, the prediction, obtained from the truncated kinematic relations \eqref{eq:kin} (displacements and modal velocities) using the values from the initial step, yields the following result:
\begin{align*}
    \hat{z}_1^{(1)}       & = z_1^{(0)} + \Delta t \, \dot{z}_1^{(0)} + 0.25 \Delta t^2 \ddot{z}_1^{(0)} = 0.25 \cdot (10^{-4})^2 \cdot 15747 = 0.39368 \cdot 10^{-4} \\
    \hat{\dot{z}}_1^{(1)} & = \dot{z}_1^{(0)} + 0.5 \Delta t \ddot{z}_1^{(0)} = 0.5 \cdot (10^{-4}) \cdot 15747 = 0.78737                                             \\
    \hat{z}_2^{(1)}       & = z_2^{(0)} + \Delta t \, \dot{z}_2^{(0)} + 0.25 \Delta t^2 \ddot{z}_2^{(0)} = 0.25 \cdot (10^{-4})^2 \cdot 22723 = 0.56808 \cdot 10^{-4} \\
    \hat{\dot{z}}_2^{(1)} & = \dot{z}_2^{(0)} + 0.5 \Delta t \ddot{z}_2^{(0)} = 0.5 \cdot (10^{-4}) \cdot 22723 = 1.1362
\end{align*}

This predictions $\hat{z}_i^{(1)}$ and $\hat{\dot{z}}_i^{(1)}$ are then substituted into the dynamic equations \eqref{eq:dyn1} and kinematic equations \eqref{eq:kin}, to yield the following implicit relations:
\begin{equation*}
    \begin{aligned}
        \ddot{z}_1^{(1)} & = - (8105.7)^2 z_1^{(1)} = - (8105.7)^2 (\hat{z}_1^{(1)}+ 0.25 \Delta t^2\ddot{z}_1^{(1)} ) \\
        \ddot{z}_2^{(1)} & = -(22866.2)^2 z_2^{(1)} = - (22866.2)^2( \hat{z}_2^{(1)}+ 0.5 \Delta t \ddot{z}_2^{(1)}  )
    \end{aligned}
\end{equation*}
Solving for the accelerations gives:
\begin{equation*}
    \begin{aligned}
        \ddot{z}_1^{(1)} & = \left[- (8105.7)^2 \hat{z}_1^{(1)} \right] \big/ \left[ 1 + 0.25 \cdot (8105.7)^2 \cdot \Delta t^2 \right]              \\
                         & = \left[- (8105.7)^2 \cdot 0.39368 \cdot 10^{-4} \right] \big/ \left[ 1 + 0.25 \cdot (8105.7)^2 \cdot (10^{-4})^2 \right] \\
                         & = -2221.6516
    \end{aligned}
\end{equation*}

\begin{equation*}
    \begin{aligned}
        \ddot{z}_2^{(1)} & = \left[- (22866.2)^2 \hat{z}_2^{(1)} \right] \big/ \left[ 1 + 0.25 \cdot (22866.2)^2 \cdot \Delta t^2 \right]                \\
                         & = \left[  - (22866.2)^2 \cdot 0.56808 \cdot 10^{-4} \right] \big/ \left[ 1 + 0.25 \cdot (22866.2)^2 \cdot (10^{-4})^2 \right] \\
                         & = -12874.1985
    \end{aligned}
\end{equation*}
It is then possible to correct the kinematic predictions $\hat{z}_i^{(1)}$ and $\hat{\dot{z}}_i^{(1)}$ by incorporating the above values to complete the equations and obtain the modal displacements and velocities at the first time step:
\begin{equation*}
    \begin{aligned}
        z_1^{(1)}       & = \hat{z}_1^{(1)} + 0.25 \, \Delta t^2 \, \ddot{z}_1^{(1)}
        = 0.39368 \cdot 10^{-4} + 0.25 \cdot (10^{-4})^2 \cdot (-2221.6516)
        = 3.3813 \cdot 10^{-5}                                                          \\
        \dot{z}_1^{(1)} & = \hat{\dot{z}}_1^{(1)} + 0.5 \, \Delta t \, \ddot{z}_1^{(1)}
        = 0.78737 + 0.5 \cdot 10^{-4} \cdot (-2221.6516)
        = 0.6763                                                                        \\
        z_2^{(1)}       & = \hat{z}_2^{(1)} + 0.25 \, \Delta t^2 \, \ddot{z}_2^{(1)}
        = 0.56808 \cdot 10^{-4} + 0.25 \cdot (10^{-4})^2 \cdot (-12874.1985)
        = 0.49245 \cdot 10^{-4}                                                         \\
        \dot{z}_2^{(1)} & = \hat{\dot{z}}_2^{(1)} + 0.5 \, \Delta t \, \ddot{z}_2^{(1)}
        = 1.1362 + 0.5 \cdot 10^{-4} \cdot (-12874.1985)
        = 0.4925
    \end{aligned}
\end{equation*}
\end{document}